\begin{abstract}
Reference-based annotation transfer via orthology has been highly successful for
protein-coding genes but remains challenging for non-coding RNAs (ncRNAs) due to
weak sequence conservation, ambiguous transcript boundaries, and incomplete
annotations. Here we present \curia{}, a pipeline for cross-species ncRNA
annotation that combines genome alignment chains with embedding-based
representations derived from RNA foundation models. \curia{} formulates
cross-species ncRNA correspondence as synteny-constrained matching in embedding
space. It first identifies candidate loci using a simplified TOGA-based
projection, then detects localized subregions (``islands'') within transcripts,
and finally establishes cross-species correspondence using similarity of
embedding representations. Because the embedding model is largely stable
across sequence context, each island is embedded once and matched at nucleotide
resolution. Using a conservative annotation-based evaluation across 19
mammalian genomes, \curia{} identified short ncRNA loci; in the human-to-mouse
comparison, 52.4\% of predicted loci overlapped an annotated exon (any overlap).
It further uncovered ${\sim}230$ candidate intergenic loci potentially
corresponding to conserved ncRNA elements absent from current annotations. These
results suggest that combining synteny with embedding-based similarity can
complement conventional sequence conservation for cross-species ncRNA annotation
in settings characterized by weak or heterogeneous primary-sequence conservation,
establishing a proof of concept for integrating comparative genomics with RNA
foundation-model embeddings for genome-scale ncRNA analysis.
\end{abstract}
